\section*{Chapitre \ref{ch:g12:ka-and-pka} --- Acides et bases~: Ka et pKa}

\begin{solution}{exo:g12:ka-and-pka:1}
Bases conjuguées~: \ce{HCOO-}, \ce{NH3}, \ce{OH-}, \ce{CO3^{2-}}.
Acides conjugués~: \ce{NH4+}, \ce{H2O}, \ce{H3O+}, \ce{CO2,H2O}.
\omterm{def:g12:ka-and-pka:oxonium}{Ampholytes}~: \ce{H2O}
et \ce{HCO3-}.
\end{solution}

\begin{solution}{exo:g12:ka-and-pka:2}
\omterm{def:g9:acids-bases-ph:ph}{pH} 1.0~; 3.0~; $-\log(\num{5.0e-3}) = 2.30$.
\end{solution}

\begin{solution}{exo:g12:ka-and-pka:3}
$10^{-4.5} = \qty{3.2e-5}{mol/L}$. À \omterm{def:g9:acids-bases-ph:ph}{pH} 11.2~:
$[\ce{H3O+}] = \qty{6.3e-12}{mol/L}$ et
$[\ce{OH-}] = \num{1.0e-14}/\num{6.3e-12} = \qty{1.6e-3}{mol/L}$.
\end{solution}

\begin{solution}{exo:g12:ka-and-pka:4}
\ce{HCOOH + H2O <=> HCOO- + H3O+},
$K_a = \dfrac{[\ce{HCOO-}][\ce{H3O+}]}{[\ce{HCOOH}]\,c^\circ}$~;
\ce{NH4+ + H2O <=> NH3 + H3O+},
$K_a = \dfrac{[\ce{NH3}][\ce{H3O+}]}{[\ce{NH4+}]\,c^\circ}$.
\end{solution}

\begin{solution}{exo:g12:ka-and-pka:5}
L'acide méthanoïque (plus petit $pK_a$). La base la plus forte est
\ce{CH3COO-}, la base de l'acide le plus faible.
\end{solution}

\begin{solution}{exo:g12:ka-and-pka:6}
$[\ce{A-}] = [\ce{H3O+}] = \qty{1.0e-3}{mol/L}$~;
$[\ce{HA}] = 0.050 - 0.001 = \qty{0.049}{mol/L}$~;
$K_a = \num{1.0e-6}/0.049 = \num{2.0e-5}$~; $pK_a = 4.69$.
\end{solution}

\begin{solution}{exo:g12:ka-and-pka:7}
$K_a = \num{6.31e-5}$~; $x^2 + \num{6.31e-5}\,x - \num{1.26e-6} = 0$
donne
$x = \qty{1.09e-3}{mol/L}$~: \omterm{def:g9:acids-bases-ph:ph}{pH} 2.96.
\end{solution}

\begin{solution}{exo:g12:ka-and-pka:8}
$\mathrm{pH} = 14.0 + \log(\num{2.0e-3}) = 14.0 - 2.70 = 11.3$.
\end{solution}

\begin{solution}{exo:g12:ka-and-pka:9}
\ce{HF}, \ce{HCOOH}, acide ascorbique, \ce{C6H5COOH}, \ce{CH3COOH},
\ce{CO2,H2O}, \ce{NH4+}, \ce{HCO3-}. Un \omterm{def:g12:ka-and-pka:strong-acid}{acide fort} se place tout en bas,
au-delà de l'extrémité de l'échelle~: il réagit totalement avec l'eau.
\end{solution}

\begin{solution}{exo:g12:ka-and-pka:10}
$[\ce{H3O+}] = 10^{-5.6} = \qty{2.5e-6}{mol/L}$, soit 25 fois plus que
dans l'eau pure (\qty{1.0e-7}{mol/L}).
\end{solution}

\begin{solution}{exo:g12:ka-and-pka:11}
$K_a K_b = \dfrac{[\ce{NH3}][\ce{H3O+}]}{[\ce{NH4+}]c^\circ} \times
\dfrac{[\ce{NH4+}][\ce{OH-}]}{[\ce{NH3}]c^\circ} = K_e$, donc
$pK_a = pK_e - pK_b = 14.0 - 4.74 = 9.26$.
\end{solution}

\begin{solution}{exo:g12:ka-and-pka:12}
À \qty{0.0010}{mol/L}~: $x^2 + \num{1.74e-5}\,x - \num{1.74e-8} = 0$,
$x = \qty{1.24e-4}{mol/L}$, \omterm{def:g9:acids-bases-ph:ph}{pH} 3.91~: une hausse de 0.52, et non de 1.
Dilué, un \omterm{def:g12:ka-and-pka:strong-acid}{acide faible} réagit en plus grande proportion avec l'eau (ici
\qty{12}{\percent} contre \qty{4}{\percent}), ce qui compense en partie
la \omterm{def:g10:concentration-and-dilution:dilution}{dilution}.
\end{solution}

\begin{solution}{exo:g12:ka-and-pka:13}
En tant qu'acide~: \ce{HCO3- + H2O <=> CO3^{2-} + H3O+}, $K = 10^{-10.33}$. En tant que
base~: \ce{HCO3- + H2O <=> H2CO3 + OH-} (\ce{H2CO3} désignant \ce{CO2,H2O}), $K = K_e/K_a = 10^{-14.0+6.37}
= 10^{-7.63}$. La réaction en tant que base a la plus grande
constante~: la solution est légèrement basique.
\end{solution}

\begin{solution}{exo:g12:ka-and-pka:14}
Un acide, aussi dilué soit-il, ne peut pas rendre une \omterm{def:g9:acids-bases-ph:acidic}{solution basique}.
Quand les \omterm{def:g12:ka-and-pka:oxonium}{ions oxonium} apportés par l'acide deviennent moins nombreux
que les \qty{1e-7}{mol/L} fournis par l'eau, les \omterm{def:g9:ions:ion}{ions} de l'eau
l'emportent~: le \omterm{def:g9:acids-bases-ph:ph}{pH} s'approche de 7 par valeurs inférieures et ne le
dépasse jamais.
\end{solution}

\begin{solution}{exo:g12:ka-and-pka:15}
$n = 0.500 / 176.0 = \qty{2.84e-3}{mol}$, $c = \qty{0.0142}{mol/L}$~;
$K_a = \num{9.12e-5}$~; $x^2 + \num{9.12e-5}\,x - \num{1.30e-6} = 0$
donne
$x = \qty{1.09e-3}{mol/L}$~: \omterm{def:g9:acids-bases-ph:ph}{pH} 2.96.
\end{solution}

\begin{solution}{pb:g12:ka-and-pka:1}
\textbf{1.} \ce{H-C(=O)-O-H}~: l'hydrogène cédé est celui de \ce{O-H}.

\textbf{2.} \ce{HCOOH}/\ce{HCOO-}~; \ce{HCOOH + H2O <=> HCOO- + H3O+}.

\textbf{3.} $K_a = \dfrac{[\ce{HCOO-}][\ce{H3O+}]}{[\ce{HCOOH}]c^\circ} =
10^{-3.74} = \num{1.82e-4}$.

\textbf{4.} En dessous de l'acide éthanoïque ($3.74 < 4.76$)~: plus fort.

\textbf{5.} $M = \qty{46.0}{g/mol}$~: $0.010 \times 46.0 = \qty{0.46}{g}$.

\textbf{6.} \ce{HCOOH}~: $0.010 - x$~; \ce{HCOO-}~: $x$~; \ce{H3O+}~: $x$.

\textbf{7.} $\dfrac{x^2}{0.010 - x} = \num{1.82e-4}$.

\textbf{8.} $x^2 + \num{1.82e-4}\,x - \num{1.82e-6} = 0$~:
$x = \qty{1.26e-3}{mol/L}$.

\textbf{9.} $\mathrm{pH} = -\log(\num{1.26e-3}) = 2.90$.

\textbf{10.} $\num{1.26e-3}/0.010 = \qty{12.6}{\percent}$.

\textbf{11.} \qty{1.26e-3}{mol/L}, c'est plus de dix mille fois
\qty{1e-7}{mol/L}~: les \omterm{def:g9:ions:ion}{ions} propres à l'eau sont négligeables.

\textbf{12.} \omterm{def:g9:acids-bases-ph:ph}{pH} 2.0.

\textbf{13.} $0.010 / \num{1.26e-3} \approx 8$ fois plus.

\textbf{14.} L'acide chlorhydrique est un \omterm{def:g12:ka-and-pka:strong-acid}{acide fort}, entièrement
transformé en \omterm{def:g9:ions:ion}{ions}~; l'acide méthanoïque est faible~: sa réaction avec
l'eau s'arrête à un équilibre où il reste en majeure partie sous forme
de \ce{HCOOH}.

\textbf{15.} \ce{HCOOH + HCO3- -> HCOO- + CO2 + H2O}.

\textbf{16.} En multipliant numérateur et dénominateur de
\[
  K = \frac{[\ce{HCOO-}]\,[\ce{CO2}]}{[\ce{HCOOH}]\,[\ce{HCO3-}]}
\]
par $[\ce{H3O+}]$, on obtient $K = K_a(\ce{HCOOH}) / K_a(\ce{CO2,H2O}) =
\num{1.82e-4}/\num{4.3e-7} \approx 420$.

\textbf{17.} Oui, $K \gg 1$~; du dioxyde de carbone se dégage en bulles.

\textbf{18.} L'\omterm{def:g9:ions:ion}{ion} hydroxyde est une \omterm{def:g12:ka-and-pka:strong-acid}{base forte} qui attaquerait
elle-même la peau~; l'\omterm{def:g9:ions:ion}{ion} hydrogénocarbonate est une \omterm{def:g12:ka-and-pka:strong-acid}{base faible} qui
neutralise l'acide et reste douce même en excès.

\textbf{19.} \omterm{def:g9:acids-bases-ph:ph}{pH} 2.9.
\end{solution}
